% The moduli of the sieve: the table of a prime and the loop of the tables, prime powers and
% products as moduli, the cost model in doubles, the choice of the moduli and of the cut points of
% the stages, the blocks of the run and the correction of the moduli between them.
\chapter{The choice of the moduli and the blocks}\label{chap:moduli}

The sieve of one denominator (Chapter~\ref{chap:sieve}) removes a numerator when the row of one of
its moduli does not admit it. This chapter says what the moduli are, how the search chooses them, as
ratpoints does in \verb|sieving_info|, and how it runs over the denominators in blocks, between
which it may correct its choice, as \verb|adapt_primes| does. A modulus is a prime of the table of
ratpoints, a power of such a prime or a product of such powers; its rows come from tables computed
from the curve (Chapter~\ref{chap:rows}). The search ranks the candidates by a model of what each
costs and what it says, computed in doubles.

None of this changes the points that the search finds or their order, whatever the values of the
doubles: the theorems hold for every choice that the search makes (Section~\ref{sec:moduliOutput}).

\section{The table of a prime and the loop of the tables}\label{sec:moduliPrimeTable}

Write $f = \sum_{i \le d} c_i x^i$ and $F(x, z) = \sum_{i \le d} c_i x^i z^{D - i}$ with
$D = d + (d \bmod 2)$ (Definition~\ref{def:binaryForm}). For a point $(a, y, b)$ with
$\gcd(a, b) = 1$ and a prime $p \nmid b$, $F(a, b) \equiv b^D f(a b^{-1}) \bmod p$, and~$b^D$ is the
square of a unit, so that $f(a b^{-1})$ is a square modulo~$p$, zero included; for $p \mid b$,
$p \nmid a$. These two facts are all that the rows of a prime rest on.

\begin{definition}[The table of a prime]\label{def:tableOfPrime}
  \lean{Ratpoints.primeTable, Ratpoints.maxBitsInPrime, Ratpoints.numPrimes, Ratpoints.Sieve.Entry, Ratpoints.Sieve.examinePrime, Ratpoints.Sieve.Entry.admits}
  \leanok
  The primes of the sieve are the odd primes below~$1024$ in increasing order
  (\csource{primes.h}{29}{40}), which \verb|decide| checks; ratpoints chooses its moduli among the
  $\mathrm{numPrimes}$ of them below $2^{\mathrm{maxBitsInPrime}}$, the~$53$ primes up to~$251$ for
  the~$8$ bits of the default build, whose Makefile passes \verb|PRIME_SIZE| $= 8$ in place of the
  header's own default~$7$ (\csource{rp-private.h}{321}{323}, \csource{rp-private.h}{325}{361}).
  The table of a prime~$p$ has $p + 1$ entries: entry $x < p$ says whether $f(x)$ is a square
  modulo~$p$, and entry~$p$ whether there are points at infinity modulo~$p$, that is, whether~$d$
  is odd or~$c_d$ is a square modulo~$p$; it is kept with~$p$ and the inverses modulo~$p$
  (\csource{rp-private.h}{542}{555}). The row of the denominator~$b$ admits the numerator~$a$ when,
  for $r = b \bmod p$ and $x = a \bmod p$, the entry at $x r^{-1} \bmod p$ is true if $r \ne 0$,
  and $x \ne 0$ if $r = 0$. This admission test is what the bit of~$a$ in the row of~$r$ says in
  ratpoints.
\end{definition}

ratpoints computes this table in \verb|examine_prime|, by Horner's rule with Barrett's reduction;
that it is the table of Definition~\ref{def:tableOfPrime} is shown in Chapter~\ref{chap:rows}.

\begin{definition}[The loop of the tables]\label{def:tablesLoop}
  \uses{def:tableOfPrime, def:sieveEntry, def:costModel, def:IsPoint, def:Options, def:bitArrays, def:forbidden, def:reciprocal, def:rowStore, def:runTables, lem:exitByPrime, lem:hornerBarrett, lem:primeRows}
  \lean{Ratpoints.TableLoop, Ratpoints.examineInLoop, Ratpoints.raiseLimit, Ratpoints.tablesLoop, Ratpoints.TableLoop.WF, Ratpoints.tablesLoop_spec}
  \leanok
  The loop of \verb|sieving_info| over the primes of the table while fewer than the number to look
  at, the normalized value of \verb|-p|, have been looked at (\csource{find_points.c}{2601}{2605}).
  The examination of a prime (\csource{find_points.c}{2606}{2626}) gives its table, the coefficients
  modulo~$p$, its density (the share of the numerators that its rows let through) and a code: $-1$,
  the exit, when the curve has no points modulo~$p$ (Chapter~\ref{chap:setup}); $0$ when the prime
  says nothing; $2$ when it says only that numerator and denominator are coprime; $1$ otherwise. A
  prime that says something becomes a candidate with its key for the first stage, and its table of
  rows goes into the store; while fewer than their bound have been found, the coefficients
  modulo~$p$ give the forbidden divisor of~$p$ (Chapter~\ref{chap:denominators}). When the caller
  left the number of primes to the library, the loop looks at one more prime, at most
  $\mathrm{numPrimes}$ in all, while the first two stages would be short of candidates that say
  more than coprimality (\csource{find_points.c}{2694}{2705}).

  The contract: for options that look at no more than $\mathrm{numPrimes}$ primes, an exit means
  that the curve has no points of any height; otherwise at most $\mathrm{numPrimes}$ primes are to
  be looked at, each table that the loop keeps, that of a prime that says something, is that of
  Definition~\ref{def:tableOfPrime} for one of the first $\mathrm{numPrimes}$ primes, every
  candidate is a prime of the table with the sieve entry of its table, the forbidden divisors are
  sound, and the store holds the tables of rows of the primes.
\end{definition}

\section{Prime powers and products}\label{sec:moduliComposite}

ratpoints also sieves with odd composite moduli from~$9$ up to $\mathrm{compositeMax} = 64$,
products of at most three powers of primes looked at. The tables of a power $q = p^e$ of a prime,
which \verb|examine_power| computes (Chapter~\ref{chap:rows}), test the numerator against~$f$ when
$p \nmid b$, and the denominator against the reversed polynomial
$\tilde f(t) = \sum_k c_k t^{D - k}$ when $p \mid b$, since then $p \nmid a$ and
$F(a, b) \equiv a^D \tilde f(b a^{-1}) \bmod q$. No Chinese remainder theorem is needed for a
product: its entry only intersects the conditions of its factors, and each factor needs only that
its modulus divides the product.

\begin{definition}[The sieve entry of a modulus]\label{def:sieveEntry}
  \uses{def:tableOfPrime, def:bitArrays, def:reciprocal, def:rowStore}
  \lean{Ratpoints.Sieve.Power, Ratpoints.Sieve.Power.admits, Ratpoints.Sieve.bitAt, Ratpoints.Sieve.Factor, Ratpoints.Sieve.SieveEntry, Ratpoints.Sieve.SieveEntry.admits, Ratpoints.Sieve.survives, Ratpoints.primeEntry}
  \leanok
  The tables of a prime power $q = p^e$ are the bit arrays \verb|fsq|, bit $x < q$ set when $f(x)$
  is a square modulo~$q$, and \verb|gsq|, bit~$t$ set when $p \mid t$ and $\tilde f(t)$ is a square
  modulo~$q$, with the inverses of the units (\csource{rp-private.h}{531}{534}). Its row of~$b$
  admits~$a$, for $r = b \bmod q$ and $x = a \bmod q$, when bit $x r^{-1}$ of \verb|fsq| is set if
  $p \nmid r$, and when $p \nmid x$ and bit $r x^{-1}$ of \verb|gsq| is set if $p \mid r$. A
  factor is a prime with its table or a prime power with its tables. The sieve entry of a
  modulus~$m$ holds its factors, the reciprocal of~$m$, with which the sieve reduces, the fields of
  the packing (Chapter~\ref{chap:sieve}), the reference to its table of rows and the kind of its
  builder (Chapter~\ref{chap:rows}), its density, and the count of its factors, $0$ for a prime,
  $1$ for a prime power and their number for a product (\csource{rp-private.h}{542}{555}). Its row
  of~$b$ admits~$a$ when the row of every factor admits $a \bmod m$ at $b \bmod m$; a numerator
  survives a list of entries for~$b$ when the row of~$b$ of every entry admits it, the interface
  between the sieve and the check.
\end{definition}

The composite moduli that the choice offers (\lname{Ratpoints.addModuli}) are the odd~$m$ from~$9$
up to $\mathrm{compositeMax} = 64$ and below $2^{\mathrm{maxBitsInPrime}}$ whose table the run can
pay for and that factor into at most three powers of primes looked at, each among the first~$64$
primes of the table and each saying something (\csource{find_points.c}{2346}{2349}); the entry of
one that the choice takes is made from the codes of its factors (\lname{Ratpoints.makeModulus}),
with the inverse of $\mathrm{rbaLength}$ modulo~$m$ by the extended Euclidean algorithm
(\csource{find_points.c}{2414}{2426}, \csource{find_points.c}{2434}{2470}).

\begin{lemma}[The rows of a modulus let every point through]\label{lem:sieveEntry_sound}
  \uses{def:tableOfPrime, def:sieveEntry, def:binaryForm, def:bitArrays, def:reciprocal, def:rowStore, lem:modinv}
  \lean{Ratpoints.Sieve.admits_examinePrime_of_sq_eq, Ratpoints.isSquare_sum_mul_pow_of_sq_eq_binaryForm, Ratpoints.Sieve.admits_run_examinePower_of_sq_eq, Ratpoints.Sieve.SieveEntry.IsSound, Ratpoints.Sieve.Factor.IsSound, Ratpoints.Sieve.isSound_prime_examinePrime, Ratpoints.Sieve.isSound_power_run_examinePower, Ratpoints.IsModulusEntry, Ratpoints.isModulusEntry_of_forall_isSound, Ratpoints.IsModulusEntry.isSound}
  \leanok
  Let~$b$ be a natural number and~$a$, $y$ integers with $\gcd(a, b) = 1$ and $y^2 = F(a, b)$.
  \begin{enumerate}
    \item For a prime~$p$, the row of~$b$ in the table of~$p$ admits~$a$.
    \item For a prime~$p$ and $e \ge 1$ with~$p^e$ within the words of the bit arrays, the row
      of~$b$ in the tables of~$p^e$ admits~$a$.
    \item A factor is sound for~$m$ when its row of $b \bmod m$ admits $a \bmod m$ for every such
      $(a, y, b)$, its inverses serve the scaling of the rows by a unit modulo~$m$, its builder has
      what it needs, and its tables are those of Definition~\ref{def:tableOfPrime} or of
      \verb|examine_power|. A prime of the sieve that divides~$m$ is a sound factor for~$m$, and so
      is a power $p^e \mid m$ of one, for $e \ge 1$ and $0 < m < 2^{\mathrm{maxBitsInPrime}}$.
    \item A sieve entry is sound when its modulus is positive and its row of~$b$ admits~$a$ for
      every such $(a, y, b)$. An entry of an odd modulus $1 < m < 2^{\mathrm{maxBitsInPrime}}$
      with sound factors whose moduli multiply to~$m$, the reciprocal and the bias of~$m$ and an
      inverse $x < m$ of $\mathrm{rbaLength}$ modulo~$m$ is an entry of~$m$ as the choice makes it
      ($\mathrm{IsModulusEntry}$), whatever its density, count, reference and builder; such an
      entry is sound.
  \end{enumerate}
\end{lemma}

\begin{proof}
  \uses{def:tableOfPrime, def:sieveEntry, lem:productRows}
  \leanok
  The entry of the table of~$p$ at $t < p$ is true whenever $f(t)$ is a square modulo~$p$. For
  $p \nmid b$, $b$ is a unit modulo~$p^e$, and in any commutative ring $F(k b, b) = b^D f(k)$, with
  $b^D$ the square of the unit~$b^{D/2}$; so $f(a b^{-1})$ is a square. For $p \mid b$,
  $p \nmid a$, which is all that the table of~$p$ asks; modulo~$p^e$,
  $F(a, b) = a^D \tilde f(b a^{-1})$, by the same identity with~$a$ and~$b$ exchanged. Zero counts
  as a square, so a high power of~$p$ dividing $F(a, b)$ needs no case of its own. A factor whose
  modulus divides~$m$ reads $b \bmod m$ and $a \bmod m$ as it reads~$b$ and~$a$, and an entry
  admits~$a$ when every factor does.
\end{proof}

\section{The cost model}\label{sec:moduliCosts}

ratpoints ranks the candidate moduli by what each costs per word of numerators swept, divided by
what it says: a modulus of density~$r$ says $-\log r$.

\begin{definition}[The cost model]\label{def:costModel}
  \uses{def:sieveEntry, def:Double, def:bitArrays, def:packing}
  \lean{Ratpoints.ModulusCandidate, Ratpoints.primeCost, Ratpoints.primeKey, Ratpoints.phase1Key, Ratpoints.phase1Wants, Ratpoints.phase2Offset, Ratpoints.checkCost, Ratpoints.runShape}
  \leanok
  A candidate modulus carries its density, key and cost, its mask of primes, its sieve entry and
  the codes of its factors (\csource{find_points.c}{147}{149}). Its cost per word swept, in units of
  an \verb|and| of the first stage, is a part for every word, the cost of its entry in the list of
  the moduli of a denominator and, in the first two stages, the costs of its set-up and, for a
  positive cost of a row, of its table of rows, spread over the~$U$ words and the~$V$ denominators
  of the run (\csource{find_points.c}{1220}{1235}); its key is that cost over $-\log r$
  (\csource{find_points.c}{1241}{1247}), with a variant for the first stage
  (\csource{find_points.c}{1256}{1270}). The rule of the first stage always takes the first
  candidate and then the next while the survivors per word that it removes, weighted by what a
  surviving word costs further on, exceed a target times its cost
  (\csource{find_points.c}{1583}{1589}); the second stage adds as many moduli as an offset, $11$ by
  default, scaled by $1/(1 + u_0/U)$ (\csource{find_points.c}{1619}{1626}). The run is estimated
  from~$64$ samples of each segment of the loop over the denominators
  (\csource{find_points.c}{1939}{1944}), and the cost of an exact check from the bit lengths of the
  coefficients and of~$H$ (\csource{find_points.c}{1285}{1318}). Every expression is computed in
  \verb|Float|, in the order of the operations of the C code; there is no contract.
\end{definition}

\begin{remark}[Why no proof reads a double]\label{rem:moduliFloats}
  The doubles are values of Lean core's \verb|Float|, IEEE~754 binary64 with the arithmetic of the
  platform, as the C code, compiled without fused multiply-adds and without \verb|-ffast-math|,
  computes them; \verb|log| and \verb|pow| are the functions of the platform's C library. No
  contract of the choice, the blocks or the correction reads the value of a \verb|Float|: each
  states only what the natural numbers of the C code do, the cut points staying within the pool and
  the entries from the end of the second stage on counting no factor, which holds whatever the
  doubles are, NaN included. So the theorems hold for every choice that the cost model may make, and
  nothing is proved about the quality of the choice. Since glibc's \verb|log| and \verb|pow| are
  not correctly rounded and glibc chooses their variant by the processor, the choice of ratpoints
  depends on the machine; the model calls the same routines of the same library (on glibc an older
  symbol version of them, with the same values on every argument that occurs here), and the
  differential tests compare the choices of the two programs on one machine
  (Chapter~\ref{chap:ccode}).
\end{remark}

\section{The choice of the moduli and of the cut points}\label{sec:moduliChoice}

In \verb|sieving_info| (\csource{find_points.c}{2522}{3103}) the first stage takes candidates in
the order of their keys, and so does the second, after the rest is keyed again by the cost of a word
that survives the first; each drops every candidate that shares a prime with one taken, so that the
moduli of the first two stages are pairwise coprime. The third stage takes the prime of least
density left while each pays for itself; the primes left after it are the reserve of the correction
(Section~\ref{sec:moduliCorrection}).

\begin{definition}[The choice of the moduli]\label{def:choiceOfModuli}
  \uses{def:Options, def:tablesLoop, def:sieveEntry, def:costModel, def:rowStore, def:forbidden, def:runTables, lem:hornerBarrett}
  \lean{Ratpoints.Params, Ratpoints.SievingInfo, Ratpoints.SieveOptions.tuning, Ratpoints.addModuli, Ratpoints.makeModulus, Ratpoints.takeEntries, Ratpoints.firstStages, Ratpoints.thirdStage, Ratpoints.choiceEnd, Ratpoints.choiceTail, Ratpoints.sievingInfo}
  \leanok
  The choice for the curve, with the options of the sieve as the search normalizes them
  (Definition~\ref{def:Options}), takes these steps:
  \begin{enumerate}
    \item the tuning, from \verb|-r|, \verb|-R|, \verb|-U| and \verb|-C| or the defaults
      (\csource{find_points.c}{2569}{2575});
    \item the first estimate, the loop of the tables with its exit, which returns the prime, and
      the further forbidden divisors (\csource{find_points.c}{2522}{2532});
    \item the estimate and the keys again, the composite moduli with their entries (\verb|add_moduli|,
      Section~\ref{sec:moduliComposite}) and the sort
      (\csource{find_points.c}{2808}{2812}, \csource{find_points.c}{2816}{2823},
      \csource{find_points.c}{2826}{2826});
    \item the first stage, the value of \verb|-n| or as many as its rule wants, and the second, the
      value of \verb|-N| or the offset, each taking candidates in their order
      (\csource{find_points.c}{2490}{2508}, \csource{find_points.c}{2843}{2846},
      \csource{find_points.c}{2864}{2871});
    \item the third stage, up to the value of \verb|-P| more or while each prime pays for itself,
      looking at further primes, with the exit, when no prime left pays
      (\csource{find_points.c}{2973}{2975}, \csource{find_points.c}{3009}{3016});
    \item the end: the third stage and the reserve sorted by density from~$s_2$ on, the entries of
      all the candidates, their reciprocals and the store with the tables of rows of every entry
      (\csource{find_points.c}{3024}{3034}).
  \end{enumerate}
  The result holds the cut points $s_1 \le s_2 \le s_3$, the ends of the three stages, the number of
  primes looked at and the estimates of the run (\csource{ratpoints.h}{395}{397}), with the sieve
  entries of all the moduli, the reserve included, their reciprocals, the forbidden divisors, the
  store of the tables of rows and the switch of the checked loop.
\end{definition}

\begin{lemma}[The choice is what the run needs]\label{lem:choiceOfModuli}
  \uses{def:Options, def:choiceOfModuli, def:tablesLoop, lem:sieveEntry_sound, lem:runRows, def:IsPoint, def:blocks}
  \lean{Ratpoints.IsSievingInfo, Ratpoints.Sieve.EntriesWF, Ratpoints.Sieve.DecisionWF, Ratpoints.IsChoiceEntry, Ratpoints.addModuli_spec, Ratpoints.makeModulus_spec, Ratpoints.ChoicePoolWF, Ratpoints.firstStages_spec, Ratpoints.thirdStage_spec, Ratpoints.sievingInfo_spec}
  \leanok
  For options that look at no more than $\mathrm{numPrimes}$ primes, the checked loop only for even
  degree, and any values of the doubles: an exit of the choice means that the curve has no points
  of any height; otherwise the result is what the run needs:
  \begin{enumerate}
    \item every sieve entry is sound (Lemma~\ref{lem:sieveEntry_sound}), with the tables of
      inverses of its factors and the fields that the packing reads;
    \item every entry that counts no factor has a single factor, a table of its modulus
      $1 < p \le 2^{32}$ whose table of inverses gives every nonzero residue below~$p$ a nonzero
      residue below~$p$, and the reciprocal of~$p$: the shape that the third stage of the sieve
      needs;
    \item the reciprocals are those of the moduli of the entries, in their order, and the store
      fits all the entries, the reserve included;
    \item the first decision, the ends~$s_2$ and~$s_3$ from block~$0$ on, is well formed:
      $s_1 \le s_2 \le s_3 \le$ the number of the entries, and every entry from~$s_2$ on counts no
      factor;
    \item the forbidden divisors are sound (Chapter~\ref{chap:denominators}), and the checked loop
      is on only for even degree.
  \end{enumerate}
\end{lemma}

\begin{proof}
  \uses{lem:sieveEntry_sound, def:tablesLoop, lem:runRows, def:forbidden, lem:candidate, lem:exitByPrime, lem:hornerBarrett}
  \leanok
  The loop of the tables gives the state of Definition~\ref{def:tablesLoop}. The loops of
  \verb|add_moduli| keep a pool ($\mathrm{ChoicePoolWF}$) of primes with their entries and of
  composites whose codes name tables that the choice holds and whose factors multiply to the
  modulus, so that the entry made from such codes at a position is an entry of its modulus, with
  the reference and the builder of the position ($\mathrm{IsChoiceEntry}$,
  Lemma~\ref{lem:sieveEntry_sound}). The stages keep every property of the candidates that does
  not read their key and cost and leave the cut points within the pool; the third stage exits only
  at a prime modulo which the curve has no points. The doubles enter only through values that no
  conjunct reads.
\end{proof}

So the result of the choice needs no well-formedness of its own. The search normalizes the options
of the sieve before the choice (\lname{Ratpoints.SieveOptions.normalize}), which bounds the number
of primes looked at by $\mathrm{numPrimes}$, so that the theorems need no hypothesis on the options.

\section{The blocks of the run}\label{sec:moduliBlocks}

The run cuts the steps of the loop of the shape (Chapter~\ref{chap:denominators}) into about~$1024$
blocks of consecutive steps and sifts each block with the moduli of the decision of the correction
in force for it; the model sifts the blocks in turn, as ratpoints does without threads. No proof
depends on the values of the constants of the blocks and of the correction but for the positivity
of a block length that the build fixes (\lname{Ratpoints.blockLen\_pos}).

\begin{definition}[The blocks of the run]\label{def:blocks}
  \uses{def:Shape, def:sift, def:deliver, def:callback, def:counters}
  \lean{Ratpoints.Sieve.Shape.segments, Ratpoints.Sieve.BlockPlan, Ratpoints.Sieve.blockLength, Ratpoints.Sieve.Shape.blockPlan, Ratpoints.Sieve.BlockPlan.steps, Ratpoints.Sieve.Shape.blockLoop, Ratpoints.Sieve.Decision, Ratpoints.Sieve.Decision.InForce, Ratpoints.Sieve.findDecision, Ratpoints.Sieve.runBlock}
  \leanok
  The steps of a shape form segments (\csource{find_points.c}{3812}{3837}): the~$k$ of the squares
  $k^2$, for each divisor~$s$ of the divisor shape the~$k$ of $s k^2$, both within
  $[b_{\mathrm{low}}, b_{\mathrm{high}}]$, or the words of~$64$ denominators, or the denominators.
  The blocks have length
  $\lfloor \mathrm{steps}/1024 \rfloor + 1$, at least~$16$, or the length that the build fixes
  (\csource{find_points.c}{3838}{3842}), and cut each segment into consecutive pieces
  (\csource{find_points.c}{3843}{3850}). A decision of the correction is the first block that it is
  for, with the ends~$s_2$ and~$s_3$ from that block on; the decision in force for block~$i$ is the
  last whose first block is at most~$i$ (\csource{find_points.c}{3908}{3910}). Block~$i$
  (\csource{find_points.c}{3889}{3898}) hands each denominator of its steps that passes the tests
  of the shape to the sieve of one denominator, with the sieve entries cut by the decision in force
  into the first two stages, the first ending at~$s_1$, and the third; it returns when the callback
  stops the search, and reports its counters with its decision.
\end{definition}

\begin{lemma}[A block delivers the points of its denominators]\label{lem:blocks}
  \uses{def:blocks, lem:choiceOfModuli, lem:runRows, def:deliver, def:callback, def:checkNumWith, def:exactBounds, def:counters}
  \lean{Ratpoints.Sieve.BlockPlan.WF, Ratpoints.Sieve.Shape.blockPlan_spec, Ratpoints.Sieve.Shape.flatMap_blockDens, Ratpoints.Sieve.findDecision_spec, Ratpoints.Sieve.runBlock_spec}
  \leanok
  The plan of the blocks is well formed, with blocks of at least one step, and the denominators of
  the blocks, concatenated in order, are those of the shape from $b_{\mathrm{low}}$ to
  $b_{\mathrm{high}}$ in the order of its loop. When the first decision is for block~$0$, the
  decision found for block~$i$ is the one in force. Let the callback not have stopped the search,
  the pieces have at least one bit array, the sieve entries be fit for the run (of the shape that
  the third stage needs where they count no factor, with the tables of inverses of their factors,
  their reciprocals in order and the row offsets of the classes; soundness is not needed here), the
  word of the classes have bit~$j$ set exactly when class~$j$ admits a numerator, the decision in
  force be well formed, the forbidden divisors be as the shape of the words needs them, $i$ be below
  the number of the blocks, $1 \le b_{\mathrm{low}}$, $b_{\mathrm{high}} < 2^{63}$, $H < 2^{63}$,
  and the store fit all the entries. Then block~$i$ keeps the store so and delivers, in order, over
  the denominators of its steps that pass the tests of the shape, the points that the exact check
  finds among the numerators that survive the first~$s_3$ entries of its decision, for every choice
  of the words of the classes and the ranges of the numerators that the packing and the bounds
  give; and when the callback has not stopped the search within it, its report is a closed form in
  the arguments of the run, which reads neither the callback nor the flags of the output.
\end{lemma}

\begin{proof}
  \uses{def:blocks, def:sift, lem:runRows, lem:deliverList, lem:shapes}
  \leanok
  The plan is arithmetic of the quotients by the length, which is positive. The loop of each shape,
  for a step with a contract on each denominator, delivers the concatenation of their lists in
  order; a stop of the callback ends it with the delivery of a prefix, which is that of the whole
  list (Lemma~\ref{lem:deliverList}). The step is the sieve of one denominator
  (Definition~\ref{def:sift}) with the entries $E_{:s_2}$ and $(E_{s_2:})_{:s_3 - s_2}$, which
  concatenate to~$E_{:s_3}$. Building rows never removes one, so that the store keeps fitting the
  entries.
\end{proof}

\section{The correction of the moduli}\label{sec:moduliCorrection}

The sieving thread counts, for each block, the words that the first stage sweeps, the bit arrays
that survive it, the numerators extracted, those prime to the denominator, those that pass the
third stage and the calls of the sieve, all before the exact check, so that the counts depend
neither on the flags of the output nor, in a block in which the callback does not stop the search,
on the callback. From the rates that the run has shown, \verb|adapt_primes| fits the survivors per
word as a linear function of the product of the densities; with \verb|-A 2| or more (not by
default) it walks the end~$s_2$ of the second stage up through the entries beyond it, or down,
above~$s_1$ and not past a composite modulus; and without \verb|-P| it computes the end~$s_3$ of
the third stage again from the new~$s_2$. A decision made at the fold of block~$j$ is for the
blocks from $j + \mathrm{adaptLag} + 1$ on, with the lag $\mathrm{adaptLag} = 32$ of
\verb|RP_ADAPT_LAG| (\csource{find_points.c}{1379}{1381}).

\begin{definition}[The correction of the moduli]\label{def:correction}
  \uses{def:Options, def:blocks, def:choiceOfModuli, def:costModel, def:callback}
  \lean{Ratpoints.Sieve.RunCounts, Ratpoints.Sieve.adaptPrimes, Ratpoints.Sieve.Correction, Ratpoints.Sieve.runFold, Ratpoints.Sieve.sieveRun, Ratpoints.adaptLag}
  \leanok
  The counters of the run, the sums of the counters of the blocks folded
  (Definition~\ref{def:counters}), those from the extracted numerators on restarted when a block sifted with another~$s_2$ is folded
  (\csource{find_points.c}{1363}{1367}). The correction (\csource{find_points.c}{1383}{1388}), when
  the run has swept enough words, sets the next one at twice the words swept and computes the
  rates, the walk of~$s_2$ under \verb|-A 2| or more and the end of the third stage, every double
  in the order of the C code.
  The fold of block~$j$ (\csource{find_points.c}{4114}{4127}) adds its counters to the state of the
  correction and appends a decision for the blocks from $j + \mathrm{adaptLag} + 1$ on when the
  correction changes an end. The run starts from the decision of the choice for
  block~$0$, sifts each block with the decisions made before it and folds its report, returns after
  the block in which the callback stopped the search, and gives the ends~$s_2$ and~$s_3$ of the
  last block folded, the moduli used (\csource{find_points.c}{5152}{5160}).
\end{definition}

\begin{lemma}[The correction keeps the decisions well formed]\label{lem:correction}
  \uses{def:correction, lem:choiceOfModuli, lem:blocks}
  \lean{Ratpoints.Sieve.Correction.WF, Ratpoints.Sieve.adaptPrimes_spec, Ratpoints.Sieve.runFold_spec, Ratpoints.Sieve.runFold_decisions_spec, Ratpoints.Sieve.correctionAfter, Ratpoints.Sieve.blockDecision, Ratpoints.Sieve.wf_correctionAfter}
  \leanok
  For any values of the doubles: when the latest decision and the decision with $s_3 = s_2$ for the
  end~$s_2$ under which the counters of the run were counted are well formed, so is the decision
  that the correction makes. A state of the correction is well formed when its first decision is for
  block~$0$ and its decisions, the latest one and the one for the end of the counters are well
  formed; the fold of a block whose decision is one of the state keeps it so, adds one block folded
  and appends decisions only for the blocks from $\mathrm{folded} + \mathrm{adaptLag} + 1$ on. The
  state after the folds of the first~$k$ blocks, each with its report in closed form
  (Lemma~\ref{lem:blocks}), and the decision of each block are functions of the arguments of the
  run that read neither the callback nor the flags of the output, well formed when the decision of
  the choice is.
\end{lemma}

\begin{proof}
  \uses{def:correction, def:exactBounds}
  \leanok
  The walk up ends at most at the number of the entries and passes only entries from the old~$s_2$
  on; the walk down ends at least at~$s_1$ and passes only entries that count no factor; the third
  stage ends between the new~$s_2$ and the number of the entries. The contracts of these loops
  state only what the natural numbers do, also where the walks compare costs that may be NaN. The
  closed form follows by induction on the number of folds.
\end{proof}

\section{Why the choice does not change the output}\label{sec:moduliOutput}

By the block lemma (Theorem~\ref{thm:blocks}) and Lemma~\ref{lem:correction}, the run delivers,
over each denominator~$b$, the points that the exact check finds among the numerators that survive
the first $s_3(b)$ entries, with $s_3(b)$ independent of the callback and the flags
(\lname{Ratpoints.Sieve.exists\_run\_sieveRun\_eq\_deliver}); every entry is sound
(Lemma~\ref{lem:choiceOfModuli}), so every point survives, whatever the choice and the corrections.
The choice changes only which candidates reach the exact check, and so the output under
\verb|RATPOINTS_NO_CHECK|, the numbers of the moduli used and the report under \verb|-v|
(Chapter~\ref{chap:output}). The statement under \verb|RATPOINTS_NO_CHECK| that compares its output
with that of the exact check (Theorem~\ref{thm:noCheck}) compares two searches with the same
options of the sieve and the same region, which differ only in the flags of the output (the second
has the exact check on and \verb|NO_Y| set); neither the choice nor the decisions read these flags
or the callback.

With threads (\verb|-t|) ratpoints sifts the blocks apart in a pool, which the model leaves out
(Chapter~\ref{chap:ccode}). The lag and thread lemmas (Theorems~\ref{thm:lag}
and~\ref{thm:threads}) prove that blocks sifted apart, block~$i$ once the blocks before block
$i - \mathrm{adaptLag}$ are folded and each with a store of its own, and delivered in their order
give what the run delivers; that the pool of the C code works so is read off the C code, not
proved.
